
The intuition that higher-quality pre-training data yields more generalizable language models motivates substantial investment in corpus curation, but this assumption has rarely been tested under controlled conditions. This study evaluates OLMo-7B (trained on Dolma, a multi-stage curated corpus) and Pythia-6.9B (trained on The Pile, a minimally curated corpus) at matched training scale (approximately 300 billion tokens), testing whether the curated-corpus model achieves a higher MMLU/HellaSwag generalization balance ratio. Contrary to the hypothesis, Pythia-6.9B achieves the higher ratio (0.5650 vs. 0.5383; difference = $-0.0267$, 95\% bootstrap CI [$-0.0447, -0.0069]$, Cohen's d = $-2.732)$. The most structurally significant finding is that both models achieve identical HellaSwag 0-shot accuracy (0.4580) at this training scale, irrespective of corpus quality. This convergence collapses the ratio into a proxy for MMLU differences that are themselves confounded by the architectural difference between the two model families (GPT-NeoX vs. LLaMA-style). The results do not support causal attribution to corpus quality without architecture-matched controls or temporal trajectory analysis. This work characterizes the conditions under which cross-architecture, single-checkpoint comparisons fail to isolate corpus quality effects, and proposes concrete requirements for future corpus quality studies.

